\documentclass[ocean-paper.tex]{subfiles}
\begin{document}
\section{Hero Pool Size}\label{sec:multihero}
\begin{figure}
    \centering
    \begin{subfigure}[t]{0.49\linewidth}
        \vskip 0pt % this vskip is used for top-alignment
        \includegraphics[width=\textwidth]{figures/multihero_trueskill.pdf}
    \end{subfigure}
    \begin{subfigure}[t]{0.49\linewidth}
        \vskip 0pt % this vskip is used for top-alignment
        \includegraphics[width=\textwidth]{figures/multihero_speedup.pdf}
    \end{subfigure}
    \caption{\textbf{Effect of hero pool size on training speed:} \genericcaptionsentence{varying the size of the hero pool.}  
    Additional heroes slows down early training only slightly.
    The severe underperformance of the 5-hero run for the first 4k versions was not investigated in detail.
    It is likely not due to the hero count but rather some instability in that particular training run.}
    \label{fig:multihero}
\end{figure}
One of the primary limitations of our agent is its inability to play all the heroes in the game.
We compared the progress in early training from training with various numbers of heroes.
In all cases, each training game is played using an independent random sampling of five heroes from the pool for each team.
To ensure a fair comparison across the runs, evaluation games are played using only the smallest set of heroes.
Because the test environment uses only five heroes, the runs which train with fewer heroes are training closer to the test distribution, and thus can be expected to perform better; the question is how much better?

In \autoref{fig:multihero}, we see that training with more heroes causes only a modest slowdown. 
Training with 80 heroes has a speedup factor of approximately 0.8, meaning early training runs 20\% slower than with the base 17 heroes.
From this we hypothesize that an agent trained on the larger set of heroes using the full resources of compute of \cleanexpname would attain a similar high level of skill with approximately 20\% more training time.
Of course this experiment only compares the very early stages of training; it could be that the speedup factor becomes worse later in training.
\end{document}